% Generated by Listas/make_statements.py from pratica-lecture-05.md, item 13. Do not edit.
\textbf{Source.} Tadelis (2013), \emph{Game Theory: An Introduction}, Exercise 9.4 ``Online Gaming'', p.~188 {[}PDF 207{]}. Answer in the source: none (the book prints no solutions; the preface says about half are on Princeton UP's website, not checked). Reproduced verbatim.



\begin{quote}

\textbf{Online Gaming:} Consider a two-stage game between two firms that produce online games. In the first stage, they play a Cournot competition game (each chooses a quantity \(q_i\)) with demand function \(p = 100 - q\) and zero marginal production costs (\(c_i(q_i) = 0\) for \(i = 1, 2\)). In the second stage, after observing the pair \((q_1, q_2)\) and after profits have been distributed, the players play a simultaneous-move ``access'' game in which each can open its platform to allow players on the other platform to play online with players on its platform (\(O\) for player 1, \(o\) for player 2), or choose to keep its platform noncompatible (\(N\) for player 1, \(n\) for player 2), in which case each platform's players can play only with others on the same platform. If they choose \((N, n)\) then second-stage payoffs are \((0, 0)\). If only one firm chooses to open its platform, it bears a cost of \(-10\) with no benefit, since the other firm did not allow open access. Finally if both firms choose \((O, o)\) then each firm gets many more eyeballs for advertising, and payoffs for each are 2500. Both players use the same discount factor \(\delta\) to discount future payoffs.



\begin{enumerate}

\def\labelenumi{\alph{enumi}.}

\item

  Find the unique Nash equilibrium in the first-stage Cournot game and all of the pure-strategy Nash equilibria of the second-stage access game. Find all the pure-strategy subgame-perfect equilibria with extreme discounting (\(\delta = 0\)). Be precise in defining history-contingent strategies for both players.

\item

  Now let \(\delta = 1\). Find a subgame-perfect equilibrium for the two-stage game in which the players choose the monopoly (total profit--maximizing) quantities and split them equally (a symmetric equilibrium).

\item

  What is the lowest value of \(\delta\) for which the subgame-perfect equilibrium you found in (b) survives?

\item

  Now let \(\delta = 0.4\). Can you support a subgame-perfect equilibrium for the two-stage game in which the players choose the monopoly quantities and split them equally? If not, what are the highest profits that the firms can make in a symmetric equilibrium?

\end{enumerate}

\end{quote}



\textbf{Note (not part of the exercise).} With the printed payoffs the threshold in (c) is \(1/16\), so the answer to (d) at \(\delta = 0.4\) is ``yes'' and the ``if not'' branch never arises.

